% Zdroj: PDF priklady-podzim-2023.pdf, fyzická str. 19-20. Značka: specializace UI-SU. Zařazeno: S3.
\section{Evaluační metriky}
\szzotazka{podzim 2023}{31}{Evaluační metriky}

\noindent\textbf{Zadání.}\par
\begin{enumerate}
  \item Definujte základní evaluační metriky pro binární klasifikátor -- úspěšnost (accuracy),
    citlivost (sensitivity, recall), přesnost (precision), F1 skóre.
  \item Chceme vytvořit systém pro detekci podvodných plateb kreditní kartou. Máme data, kde jsou miliony
    plateb, z nichž jen několik stovek jsou platby podvodné. Jaké metriky použijete pro vyhodnocení takového
    systému? Jaké metriky naopak nebudou moc vhodné? Odpověď zdůvodněte.
  \item Pro výše zmíněnou úlohu jsme vytvořili dva systémy. Jeden z nich má citlivost $0{,}99$ a přesnost
    $0{,}2$. Druhý má citlivost $0{,}9$ a přesnost $0{,}95$. Vysvětlete, co tyto hodnoty prakticky znamenají.
    Který ze systémů byste použili? Své rozhodnutí zdůvodněte.
\end{enumerate}

\medskip
\noindent\textbf{Náčrt řešení.}\par
\begin{enumerate}
  \item Definujme TP/TN/FP/FN jako počty true/false positive/negative. Správnost potom je
    $(\mathrm{TP}+\mathrm{TN})/(\mathrm{TP}+\mathrm{FP}+\mathrm{TN}+\mathrm{FN})$, citlivost je
    $\mathrm{TP}/(\mathrm{TP}+\mathrm{FN})$, přesnost je $\mathrm{TP}/(\mathrm{TP}+\mathrm{FP})$
    a F1 skóre je harmonický průměr citlivosti a přesnosti $= 2\mathrm{TP}/(2\mathrm{TP}+\mathrm{FP}+\mathrm{FN})$.
  \item Pro vyhodnocení systému je vhodné použít dvojici citlivost a přesnost. Dá se také použít F1 skóre
    nebo AUC. Nevhodná naopak bude správnost vzhledem k výrazně nevyváženému množství instancí v jednotlivých třídách.
  \item První model odhalí $99\,\%$ podvodných plateb, na druhou stranu ale jen $20\,\%$ z ohlášených
    podezřelých plateb budou platby podvodné. Druhý model sice odhalí jen $90\,\%$ podvodných plateb,
    ale $95\,\%$ z podezřelých plateb opravdu podvody budou. V případě použití prvního modelu tedy bude
    potřeba prošetřit cca 5krát více případných podvodů, což může vést ke zbytečnému obtěžování zákazníků
    a i vyšším nákladům na prošetřování. Může tedy být vhodnější použít model druhý. Na druhou stranu,
    pokud bychom chtěli mít větší jistotu, že podvody odhalíme a výše zmíněné nevýhody nám nevadí, je
    vhodnější model první. Neexistuje tu jednoznačně správná odpověď, důležitá je především argumentace.
\end{enumerate}
